\chapter*{Permitted Use of Arabic Terms}
\addcontentsline{toc}{chapter}{Permitted Use of Arabic Terms}
\markboth{Arabic Terms in Use}{Arabic Terms in Use}

The following Arabic terms appear repeatedly throughout the book and
are used without italics after first introduction. Their meanings are
collected here for quick reference; full definitions and citations
appear in the glossary in the back matter and at the relevant
chapters.

\begin{notationbox}[Core terms]
\begin{description}[leftmargin=4cm,style=nextline]
    \item[\textbf{rizq}] provision; that which is allotted to a
        creature as material sustenance. The Quran ascribes the act of
        providing to Allah under the divine name \textit{ar-Razzaq}
        (the Provider). Plural: \textit{arzaq}. Treated formally in
        Chapter 10.
    \item[\textbf{barakah}] divine blessing; in operational terms, the
        multiplier on material inputs that distinguishes wealth (or
        time, or knowledge) of the same nominal magnitude into
        different effective magnitudes. Treated formally in Chapter 15.
    \item[\textbf{qadar}] divine decree; the determination of all
        events according to divine knowledge and will. Distinguished
        from \textit{jabr} (compulsion) by classical theology.
        Treated formally in Chapter 16.
    \item[\textbf{ikhtiyar}] human choice; the element of free will.
        Compatibly co-existing with \textit{qadar} in classical
        Islamic theology, the precise structure of compatibility being
        the subject of the Mu`tazili--Ash`ari--Maturidi disputes
        treated in Chapter 3.
    \item[\textbf{tawhid}] divine unity; the metaphysical ground of
        all things; the first article of Islamic faith. The
        philosophical articulation as \textit{wahdat al-wujud} is
        treated in Chapter 9.
    \item[\textbf{niyyah}] intention; the inner state that gives an
        act its spiritual and (in our framework) economic meaning.
        The opening hadith of \textit{Sahih al-Bukhari} establishes
        the principle that ``actions are by intention.''
    \item[\textbf{wujud}] being, existence; in Akbarian metaphysics,
        the foundational ontological category. Admits of degrees,
        with absolute \textit{wujud} (Allah) the source of all
        derived \textit{wujud}.
    \item[\textbf{wahdat al-wujud}] the unity of being; Ibn Arabi's
        central doctrine that there is, ultimately, only one
        \textit{wujud}, and that all created things are loci of its
        self-disclosure. Distinguished from pantheism in Chapter 9.
    \item[\textbf{tajalli}] theophany; divine self-disclosure; the
        active process by which absolute \textit{wujud} manifests in
        finite loci of existence. Each moment of created existence is
        a fresh \textit{tajalli}.
    \item[\textbf{insan al-kamil}] the perfect human; the locus that
        manifests the divine names in equilibrium; in Akbarian
        thought, exemplified by the Prophet Muhammad.
    \item[\textbf{riba}] interest, usury; the gratuitous return on
        debt, prohibited in Islamic law. Distinguished into
        \textit{riba al-fadl} (excess in exchange of like commodities)
        and \textit{riba al-nasi'ah} (interest on a deferred
        repayment). Treated formally in Chapter 13.
    \item[\textbf{zakat}] the obligatory annual alms in Islamic law;
        a proportional levy on wealth held for a year, redistributed
        to specified categories of recipients. Treated formally in
        Chapter 17.
    \item[\textbf{waqf}] charitable endowment; the perpetual
        alienation of property to a beneficiary purpose. Plural:
        \textit{awqaf}. Treated formally in Chapter 17.
    \item[\textbf{takaful}] mutual cooperative insurance; the
        Islamic alternative to conventional insurance, structured as
        risk-pooling rather than risk-transfer.
    \item[\textbf{mudarabah}] a profit-sharing partnership in which
        one party provides capital and the other expertise; profits
        are shared by agreed ratio, losses borne by the capital
        provider only.
    \item[\textbf{musharakah}] an equity partnership in which all
        parties contribute capital and share both profits and losses
        in proportion to their contributions.
    \item[\textbf{qard hasan}] a benevolent loan, given without
        expectation of return beyond the principal; the Quranic
        alternative to interest-bearing lending.
    \item[\textbf{maqasid al-shari`ah}] the higher objectives of the
        revealed law; classically enumerated as the protection of
        religion, life, intellect, lineage, and property.
\end{description}
\end{notationbox}

\begin{notationbox}[Other terms used regularly]
\begin{description}[leftmargin=4cm,style=nextline]
    \item[\textbf{kalam}] speculative theology; the discursive
        defense and articulation of doctrine.
    \item[\textbf{fiqh}] Islamic jurisprudence; the practical
        derivation of legal rulings from the sources.
    \item[\textbf{shari`ah}] the revealed law in its broad sense; the
        path of right action.
    \item[\textbf{`ilm}] knowledge.
    \item[\textbf{`adl}] justice.
    \item[\textbf{ihsan}] excellence; doing what is good in the most
        beautiful way.
    \item[\textbf{taqwa}] God-consciousness.
    \item[\textbf{shukr}] gratitude; the recognition of blessing and
        its proper acknowledgement.
    \item[\textbf{tawakkul}] trust in divine provision; reliance on
        Allah following the discharge of one's effort.
    \item[\textbf{sabr}] patience; perseverance through difficulty.
    \item[\textbf{ummah}] the community of believers, considered as a
        single body.
    \item[\textbf{khalifah}] vicegerent; the human role as steward of
        creation.
    \item[\textbf{amanah}] trust, in the sense of something held in
        trust; wealth, life, and station are amanat held by the
        person from Allah.
    \item[\textbf{kasb}] acquisition; in Ash`ari theology, the
        technical term for the human contribution to an action that
        is otherwise created by Allah.
    \item[\textbf{hisbah}] the institution of public-interest
        oversight in classical Islamic governance, including market
        inspection.
    \item[\textbf{muhtasib}] the official charged with hisbah;
        market inspector and inspector of public morality.
    \item[\textbf{`alam al-mithal}] the imaginal world; the
        intermediate ontological domain in Akbarian metaphysics
        between pure spirit and pure matter.
    \item[\textbf{khayal}] imagination; the faculty by which the
        imaginal world is accessed.
    \item[\textbf{barzakh}] isthmus, intermediate state; the realm
        between this life and the resurrection, and more generally
        any intermediate ontological domain.
\end{description}
\end{notationbox}

\section*{A note on plurals and word formation}

Where Arabic plurals exist and are commonly used in scholarship, both
singular and plural forms appear (\textit{rizq/arzaq},
\textit{waqf/awqaf}). Where English plurals are conventional in the
secondary literature, we use them (\textit{hadiths},
\textit{maqasid}, etc.). The reader should not infer significance
from the choice; it follows the usage of the secondary literature in
each instance.

Where Arabic terms enter English compound expressions, we follow the
convention of the secondary literature: ``zakat-al-fitr,''
``maqasid al-shari`ah.''

\section*{Capitalization}

Arabic terms are not capitalized in running text unless they appear at
the beginning of a sentence or are part of a proper name. The divine
name \textit{Allah} is capitalized throughout, as are the names of the
divine attributes (\textit{ar-Rahman}, \textit{ar-Razzaq}, etc.) when
referring to the divine attribute itself rather than to the name as a
linguistic object.

\cleardoublepage
